\chapter{Sustainability Analysis and Ethical Implications}

The sustainability profile of this thesis is defined by a tradeoff: the
exploratory computation needed to discover a reliable model versus the much
lower marginal cost of reusing it for dense CKM generation. Since the result is
a software pipeline rather than a manufactured device, the analysis focuses on
energy use, reuse of trained artifacts, reproducibility, and the risks of using
simulated predictions in planning decisions.

\section{Sustainability matrix}

\begin{table}[htbp]
\centering
\caption{Sustainability matrix summary for the thesis.}
\label{tab:sustainability_matrix}
\scriptsize
\begin{tabularx}{\textwidth}{lXXX}
\toprule
Perspective & Development of BT & Project execution & Risks and limitations \\
\midrule
Environmental &
GPU training, laptop use, storage, and writing. &
Reduced need for repeated ray tracing simulations once the surrogate is trained. &
Extra retraining, larger models, or poor reuse of cached priors could increase footprint. \\
Economic &
Student and supervision time, plus shared HPC use. &
Lower marginal planning cost if the surrogate replaces repeated simulations. &
Real deployment needs field validation and maintenance. \\
Social &
Skills gained in wireless modelling, ML, HPC workflows and scientific reporting. &
Faster elevated base station or UAV coverage planning for temporary, rural or
emergency networks. &
Bias toward represented city morphologies and risk of overclaiming simulated results. \\
\bottomrule
\end{tabularx}
\end{table}

\subsection{Environmental impact}

\textbf{Development of BT.}

The main environmental cost during development was computation. The final
\textsc{HARP-Net CKM} line used a four GPU RTX~2080~Ti node for approximately
\SI{12}{h}; the distribution first and prior calibration phase used additional
GPU/CPU time. Across the full exploratory project, the total exceeded roughly
400 GPU hours. This is a one time research overhead and includes many failed
or diagnostic experiments, not only the selected final run.

Several decisions reduced avoidable impact: reusing saved histories and frozen
calibration artifacts, evaluating city holdout metrics from stored checkpoints,
and keeping the final priors frozen instead of recomputing them blindly inside
every neural experiment. The final design is also relatively efficient at
inference. On the Barcelona benchmark, the local MATLAB ray tracing path took
\SI{99.89}{s} per full resolution sample, while the complete generator path
with masks, priors, neural prediction and export took \SI{0.88}{s}. The
sustainability advantage therefore depends on reuse: if the same calibrated
generator is queried many times, the marginal cost is much lower than repeated
ray tracing; if the model is retrained from scratch for every new scenario, that
advantage weakens.

The life cycle interpretation is therefore important. During the research phase,
the footprint is dominated by exploration: many failed or partially successful
attempts were necessary to discover that the final path should combine strong
priors with a bounded neural residual model. During the use phase, the footprint is
dominated by inference and by occasional recalibration. The project is most
sustainable when the trained model, frozen priors, cached masks and exported
checkpoints are reused across many planning queries. It would be less
sustainable if the same design were scaled by simply increasing model size or
retraining from zero for every city.

The final generator also improves reproducibility, which has an environmental
side. A documented \texttt{.npz} output, explicit mask policy, and fixed
hardware/runtime benchmark make it easier to compare changes without rerunning
expensive ray tracing baselines. Reproducible experiments are not only a
scientific requirement; they reduce wasted computation caused by ambiguous
settings, hidden preprocessing changes or repeated debugging runs.

For the development footprint, a rough quantitative estimate can be recovered
from the experiment history. The final \textsc{HARP-Net CKM} line used four
RTX~2080~Ti GPUs for approximately \SI{12}{h}, about \SI{18}{kWh} under a
\SI{1.5}{kW} four GPU node approximation. The distribution first and prior
calibration phase added roughly another \SI{18}{h} of four GPU time, about
\SI{27}{kWh}, plus lighter CPU and DirectML calibration runs. Across the full
exploratory thesis line, including failed and diagnostic experiments, the total
exceeded roughly \num{400} GPU hours. With the same power approximation this is
about \SI{150}{kWh}, or approximately \SI{30}{kgCO2e} if an illustrative
\SI{0.2}{kgCO2e/kWh} grid factor is used. This is not a full life cycle
assessment and does not include embodied carbon of shared GPUs or data centre
infrastructure, but it gives an order of magnitude for the academic research
phase.

Because this is a software thesis, raw material use is indirect: the laptop,
shared GPU node, storage and network infrastructure already existed. No dedicated
device or project specific material was purchased. Supplier environmental reports
and public codes of ethics or conduct for the shared hardware providers were not
audited, which is a limitation of the analysis. For a future operational version,
cloud, hardware or data providers should be selected with public environmental
and ethical reporting when possible.

\textbf{Project execution.}

For the use phase, the older detailed estimate remains useful because it
connects the benchmark with energy rather than only speed. Extrapolating the
Barcelona runtime to the \num{16180} CKM samples gives \SI{448.96}{h} for the
local ray tracing path and \SI{3.95}{h} for the surrogate. Under a simple
\SI{100}{W} same laptop full load assumption, this corresponds to about
\SI{44.9}{kWh} versus \SI{0.39}{kWh}. The assumptions are deliberately simple,
but the conclusion is robust in direction: the environmental case depends on
amortizing the research footprint through many reused predictions, not on
claiming that model training itself is free.

\textbf{Risks and limitations.}

At the end of life, the project mainly generates digital artefacts that can be
archived or deleted; no project specific physical waste is expected. The main
environmental risk is poor reuse. If the model is retrained from scratch for
every new city, or scaled by larger architectures without a measured need, the
environmental benefit can disappear. If the work were repeated, the footprint
could be reduced by freezing the experiment plan earlier, running smaller pilot
subsets before full training, reusing checkpoints more aggressively and
scheduling heavy jobs when cleaner electricity is available.

\subsection{Economic impact}

\textbf{Development of BT.}

Economically, this work is not yet a product ready to take to market. Its value
is as a technical prototype showing that, under the fixed \SI{7.125}{GHz} CKM
setting used here, dense ray traced targets can be approximated by a reusable
surrogate once the target city representation and calibration protocol are fixed.
The main development cost was human time:
literature review, dataset handling, training/debugging, analysis, supervision
and writing. The compute cost is small compared with this time cost in an
academic project, but it is still relevant because the full exploration
required many GPU hours.

Table~\ref{tab:sustainability_budget} gives an indicative professional cost
view of the BT development. It is not a billed project budget, but it answers
the guide's request to quantify human and material resources, including working
hours, hourly costs, computation and operating assumptions.
The approximate calendar was literature review and proposal in the first stage,
dataset handling and baseline implementation next, experiment design and
debugging in the central stage, and writing, figures and review in the final
stage, with supervision distributed throughout. No subcontracting, new computer
purchase, material purchase or separate tax item is included; electricity is
included in the shared compute estimate. Cost was reduced by using open source
software, shared university compute, cached outputs and reusable scripts instead
of buying dedicated hardware or repeating ray tracing unnecessarily.

\begin{table}[htbp]
\centering
\footnotesize
\renewcommand{\arraystretch}{1.08}
\caption{Indicative economic cost of the BT development.}
\label{tab:sustainability_budget}
\begin{tabularx}{\textwidth}{>{\raggedright\arraybackslash}Xrrr>{\raggedright\arraybackslash}p{3cm}}
\toprule
Task or resource & Quantity & Unit cost & Cost & Notes \\
\midrule
Literature review and proposal & 60\,h & 12\,EUR/h & 720\,EUR & Initial scope and state of the art \\
Dataset and baseline implementation & 90\,h & 12\,EUR/h & 1\,080\,EUR & HDF5, U-Net/cGAN, masking \\
Experiment design and debugging & 180\,h & 12\,EUR/h & 2\,160\,EUR & Priors, experts, evaluation \\
Writing and figures & 120\,h & 12\,EUR/h & 1\,440\,EUR & Thesis text and diagrams \\
Supervision meetings and supervisors' dataset simulation time & 50\,h & 40\,EUR/h & 2\,000\,EUR & Advisor time \\
GPU compute, \textsc{HARP-Net CKM} & 48 GPU hours & 0.25\,EUR/GPU hour & 12\,EUR & 4$\times$ RTX~2080~Ti, 12\,h \\
GPU compute, distribution and prior phase & 72 GPU hours & 0.25\,EUR/GPU hour & 18\,EUR & Distribution first models plus prior calibration \\
Exploratory compute & $\approx$400 GPU hours & 0.25\,EUR/GPU hour & $\approx$100\,EUR & Research iterations \\
Electricity & N/A & included in GPU hours & 0\,EUR & Covered by shared compute rate \\
Storage and code hosting & N/A & free & 0\,EUR & GitHub and local storage \\
\midrule
\textbf{Total indicative cost} & & & \textbf{$\approx$7\,530\,EUR} & \\
\bottomrule
\end{tabularx}
\end{table}

\textbf{Project execution.}

The most realistic economic path is a planning tool rather than a standalone
consumer product. A mature version could be licensed as software for network
planning teams, integrated into a digital twin workflow, or offered as a
service for rapid elevated base station or UAV coverage studies. Operators
could use faster CKM generation to compare transmitter heights above roughly
\SI{10}{m}, test temporary base station locations, identify weak streets or
squares before deployment, and prioritize the cases where full ray tracing or
field validation should be run. Better radio maps would
not create capacity by themselves, but they would make planning iterations
cheaper and could reduce the risk of sending equipment or measurement teams to
poor candidate sites. Patentability would require a novelty analysis beyond
this thesis, because CKMs, radio map prediction and hybrid physics neural
models already exist in the literature; the more plausible near term value is
know how, implementation, calibration data and integration into operator
planning workflows.

As a simplified execution budget, a first year pilot version would be dominated
by engineering and validation rather than compute. A reasonable minimum estimate
is about \SI{40}{h} of engineering packaging and documentation at 30~EUR/h
(\(\approx\)1200~EUR), \SI{10}{h} of technical validation or supervisor review
at 40~EUR/h (\(\approx\)400~EUR), less than 50~EUR of storage and routine batch
inference, and 50 to 150~EUR of shared GPU time if one recalibration or small
retraining pass is required. This gives an indicative first year maintenance
and pilot operation cost of roughly 1700 to 2000~EUR, excluding field measurement
campaigns. If real measurements are required, they become the dominant cost and
should be budgeted as a separate deployment validation task.

\textbf{Risks and limitations.}

The economic limitations are therefore not only computational. The project
would also need packaging, documentation, support, validation datasets, update
procedures and periodic recalibration if it were maintained as planning
software. The cost of dismantling is small because the project is software, but
archival responsibility remains: code, checkpoints, data provenance and model
metadata should be preserved so that future users do not have to rerun the same
experiments only to recover the assumptions.
The main viability risks are that operators may need other frequencies, antenna
patterns, map formats or field validated accuracy levels, which would increase
calibration cost. Other projects could still benefit from the reusable priors,
mask generation code, evaluation scripts and documented city holdout protocol.

\subsection{Social impact}

\textbf{Development of BT.}

From the development perspective, the work also had a personal and professional
training impact. It required propagation theory, machine learning, HPC
workflows, reproducible evaluation, scientific writing and continuous critical
review of failed approaches. Inclusive language is used where applicable, and
the thesis does not assign technical roles, expected users or beneficiaries by
gender. Since no dedicated distributors, manufacturers or suppliers were
contracted for this thesis, their public codes of ethics or conduct were not
audited; a future maintained tool should include that check when selecting
external cloud, data or hardware providers.

\textbf{Project execution.}

Socially, faster CKM prediction can support temporary coverage planning for
public events, rural extensions, emergency response or infrastructure
restoration. The direct users would be network planners; the indirect
beneficiaries would be users receiving more reliable wireless coverage.

The social value is strongest when the tool is used as decision support rather
than as an automatic decision maker. A fast surrogate can help planners compare
candidate elevated transmitter heights, identify weak coverage zones, and
prioritize detailed simulation or measurement budgets. It should not replace domain review:
the model has been validated on simulated CKM cities and inherits the
assumptions of the ray tracer, the topology representation, and the ground mask
definition.

From the project execution perspective, usability matters because the tool would
likely be used by engineers rather than by the general public: input
metadata, mask policy, antenna height, checkpoint version and limitations
should be visible enough that users with different levels of ML expertise can
understand what the generated map does and does not certify.

The current sector context is a growing use of digital twins, simulation assisted
planning and machine learning surrogates for network design. The same problem
could also be addressed with full ray tracing, measurement campaigns, empirical
propagation formulas or commercial planning tools; those options shift the
balance between cost, accuracy, transparency and accessibility. The main barriers
to this prototype are technical: users need access to map inputs, metadata and
basic interpretation of uncertainty and masks. A future interface should provide
accessible plots, documented defaults and clear warnings so that age, disability,
cultural background or limited ML expertise do not become hidden barriers.

\textbf{Risks and limitations.}

The main social risk is overconfidence: a simulated city holdout result should
not be treated as certified field performance. Errors should continue to be
reported by LoS/NLoS state, transmitter height bin and morphology group rather
than hidden behind one global RMSE. A dependency risk also exists: if planners
rely only on the surrogate and stop running measurements or high fidelity
simulation, users could become vulnerable to systematic model errors in certain
environments.

There is also a fairness dimension in the morphology split. The model is tested
on unseen cities, but the training distribution is still finite and urban
morphologies are not equally represented. Dense high rise, unusual street
canyons, informal settlements, vegetation heavy areas or non European building
patterns may behave differently from the dominant CKM regimes. For that reason,
future deployment should report uncertainty or at least split diagnostics by
morphology and height, instead of using the global average as the only quality
indicator.

\section{Ethical implications}

The exact need addressed is faster dense CKM generation for UAV and elevated base
station planning under the thesis dataset setting. This need was defined by the
BT proposal, the supervisors' research context and the broader network planning
problem of replacing repeated simulation with documented surrogates where this
is justified. Consequences beyond the intended use include misplaced confidence
in coverage maps, pressure to skip measurements, or inappropriate use outside the
frequency, city and antenna assumptions evaluated here.

The main ethical obligation is honesty about validity. The dataset is generated
by ray tracing, so the thesis demonstrates simulator level generalization to
unseen cities, not real world deployment performance. A future operational
system would require field measurements, calibration and human review,
especially before public service or emergency use.

A second ethical issue is transparency. The final predictor is not a pure
black box CNN: the attenuation and spread priors expose the geometric and
morphological assumptions before the residual network is applied. This makes the
system easier to audit, because a user can inspect whether an error comes from
the LoS/NLoS mask, the prior, or the learned residual correction. That audit
path is particularly relevant for emergency or infrastructure planning, where a
wrong map can lead to misplaced confidence in predicted coverage.

The work does not process personal data, but operational radio planning can
interact with sensitive contexts such as emergency response, public events or
critical infrastructure. The generator should therefore be accompanied by clear
metadata: input source, resolution, antenna height, mask generation mode,
checkpoint version and known limitations. Keeping that metadata with the
prediction is part of the ethical responsibility of making a fast tool: speed is
useful only if the output remains traceable.

The project also aligns with the professional ethics expected in engineering
and research work by documenting assumptions, avoiding data leakage between
cities, reporting limitations, and keeping the evaluation rules explicit.
In UPC terms, the relevant principles are those of the Code of Ethics and the
Code of Research Integrity: transparency, honesty in reporting results,
responsible use of resources, and respect for the people who may be affected by
technical decisions. A future operational version should include human review
before coverage maps affect public service, emergency response or critical
infrastructure decisions.

\section{Relationship with the Sustainable Development Goals}

The project aligns most directly with SDG~9, because it develops a faster tool
for wireless infrastructure planning; SDG~11, because it can support resilient
urban and temporary connectivity; and SDG~13, because a validated surrogate can
reduce repeated simulation energy. These contributions are indirect: the thesis
provides a reproducible method and generator, not an operational network
deployment by itself.

In the guide's language, these SDG links are indirect but specific:
SDG~9 is supported through a faster infrastructure planning method, SDG~11
through resilient urban and temporary connectivity planning, and SDG~13 through
lower repeated simulation energy when the surrogate is reused responsibly. The
work does not by itself deploy infrastructure or guarantee better coverage; it
provides a reproducible method that can help future planning tools prioritize
where slower simulation or real measurement is most needed.

Overall, the sustainability conclusion is conditional but positive. The research
process was computation heavy, but it produced a reusable generator that can
reduce many repeated full resolution simulation calls once validated for a
target domain. The strongest sustainable use case is therefore not ``train a
larger model'', but ``reuse \textsc{HARP-Net CKM} as a documented surrogate to screen
planning options and prioritize slower simulation or real measurement''.
